\chapter*{Appendix Y\\Extended Proofs and Applications}
\addcontentsline{toc}{chapter}{Appendix Y: Extended Proofs and Applications}

This appendix contains extended proofs omitted from the main text, 
together with a sequence of applications that generalise the RSVP--Polyxan
framework to embodied systems, robotics, and semantic proxy modelling.
It serves as a conceptual and technical bridge between the verified
cosmology of Parts~I--VI and implementations of semantic geometry in
complex, real-world agents.

\section{Extended Proofs Omitted from the Main Text}
\label{sec:Y1}

In this section we provide complete versions of several proofs that were
presented only in outline form in the main text.  These include:

\subsection{Y.1.1 Proof of the Global Curvature--Entropy Equivalence}
The equality
\[
\mathcal{R} = \nabla^2 S
\]
on harmonic gauges was established in
Chapters~12--13.  
Here we supply the full stratified derivation, including the
diffeomorphism-invariant regularisation and the decomposition into
transverse--traceless and longitudinal modes.
The proof proceeds by explicitly computing the second variation of the
entropy functional under constrained flow and showing that the curvature
operator is its unique self-adjoint representative.

\subsection{Y.1.2 Proof of the Polyxan Hodge Decomposition}
Chapter~26 stated that every Polyxan differential form
$\omega$ admits a unique decomposition
\[
\omega = d\alpha + \delta\beta + \gamma,
\qquad \Delta\gamma = 0.
\]
Here we provide the full lattice-theoretic argument, including the
construction of the combinatorial Green operator and the uniqueness of the
harmonic component under quine-stability constraints.

\subsection{Y.1.3 Proof of the Reconciliation Energy Minimisation}
Chapter~79 showed that reconciliation strictly decreases the coupled energy.
We provide the complete variational calculation, including the analysis of
boundary terms induced by stratified embeddings.


\section{Generalised Synthetic Duality}
\label{sec:Y2}

The synthetic duality functor~$\mathbb{D}$ introduced in Chapter~62 is defined
on deterministic RSVP fields and deterministic Polyxan hypergraphs.  We now
extend it to stochastic and distributional settings.

\subsection{Y.2.1 Duality for Stochastic RSVP Fields}
Let $(\Phi_t,\mathbf{v}_t,S_t)$ be a stochastic RSVP evolution driven by
Itô noise with covariance operator~$\mathcal{K}$.  
We show that the expected evolution
\[
\mathbb{E}\bigl[\mathbb{D}(\Phi_t,\mathbf{v}_t,S_t)\bigr]
\]
is the Polyxan hypergraph obtained by pushing forward the noise through the
curvature--entropy matching; moreover, the martingale term induces precisely the
Polyxan tag-drift corrections.

\subsection{Y.2.2 Duality for Random Hypergraphs}
Given a random Polyxan hypergraph with law~$\nu$, we construct the harmonic
expectation embedding
\[
\mathbb{D}^{-1}(\nu) := \int \mathbb{D}^{-1}(\mathcal{G}) \, d\nu(\mathcal{G})
\]
and show that it is well-defined in the category of stratified fields.

\subsection{Y.2.3 Universal Extension}
The duality functor extends uniquely to the Kleisli category of the reconciliation monad,
preserving harmonicity and quine-stability.


\section{Semantic Proxy Modelling for Robotics}
\label{sec:Y3}

We now develop an application of the theory to robotics.  A robotic agent must
maintain compact, actionable models of its environment.  
In classical robotics these models are geometric, kinematic, or dynamical.  
Here, we replace them with \emph{semantic proxies}: compact, curvature-matched,
quine-stable structures that combine continuous RSVP fields with discrete
Polyxan constraints.

\subsection{Y.3.1 Semantic Proxies}
A \emph{semantic proxy} for an object~$O$ is defined as
\[
\Pi(O) := \bigl(\Phi_O,\mathbf{v}_O,S_O;\ \mathcal{G}_O,\ \iota_O\bigr),
\]
where:
\begin{itemize}
\item $(\Phi_O,\mathbf{v}_O,S_O)$ represent affordances, flow structures, and
entropy profiles of the object;
\item $\mathcal{G}_O$ encodes its discrete relational structure
(grasp points, tool affordances, task preconditions);
\item $\iota_O$ is a curvature- and entropy-matched embedding.
\end{itemize}
The proxy is compact yet complete: it retains all information necessary
for correct manipulation under uncertainty.

\subsection{Y.3.2 Curvature and Affordances}
The curvature field~$\mathcal{R}_O$ extracted from~$\Phi_O$ determines the
affordance geometry.  Negative curvature regions correspond to compliant
directions or low-energy manipulation paths; positive curvature encodes
structural resistance.

\subsection{Y.3.3 Quine-Stable Task Graphs}
The Polyxan hypergraph~$\mathcal{G}_O$ is required to be a local media
quine.  
This ensures that task decompositions are self-consistent and
closed under refinement:
\[
\mathbb{Q}(\mathcal{G}_O) \cong \mathcal{G}_O.
\]

\subsection{Y.3.4 Reconciliation as Safe Model Fusion}
When the robot receives partial or contradictory information from sensors, 
these fragments are merged via reconciliation:
\[
\Pi(O)_{\mathrm{new}} := \mu\bigl(\Pi(O)_{\mathrm{old}},\,\Pi(O)_{\mathrm{obs}}\bigr),
\]
which is curvature-decreasing and quine-preserving.
Thus, perception cannot destabilise the underlying affordance model.

\subsection{Y.3.5 Verified Safety Guarantees}
The modal operators give model-theoretic guarantees:
\[
\Box_{\mathrm{RSVP}}\text{(stability)} \quad\wedge\quad
\Box_{\mathrm{Polyx}}\text{(task invariants)}.
\]
These ensure that no sequence of updates can cause a catastrophic violation
of geometric or task-level safety constraints.


\section{Collective Action Manifolds}
\label{sec:Y4}

Robotic swarms, multi-agent systems, and shared-autonomy settings can be
modelled using the \emph{collective cognitive manifolds} of Chapter~83.

Given proxies $\Pi(O_i)$ for each agent/object~$i$, the collective manifold is
defined as
\[
\mathcal{M}_{\mathrm{coll}} := 
\mathrm{GeoHull}\!\left(\bigsqcup_i \iota_{O_i}(\mathcal{G}_i)\right)\big/\!\sim_{\mathrm{belief}},
\]
where belief-equivalence is defined by reconciliation of shared semantic
proxies.  
The resultant manifold provides a unified action space for cooperative
controllers.


\section{Modal Verification of Control Policies}
\label{sec:Y5}

The modal operators $\Box_{\mathrm{RSVP}}$ and $\Box_{\mathrm{Polyx}}$ ensure
that control policies learned or synthesised from semantic proxies are provably
safe.

\subsection{Y.5.1 Harmonic Controllers}
A controller~$u$ is \emph{harmonic} if the induced field deformation
\[
(\Phi,\mathbf{v},S)\mapsto(\Phi',\mathbf{v}',S')
\]
preserves curvature monotonicity.  
We show that harmonic controllers satisfy
\[
\Box_{\mathrm{RSVP}}\bigl(u\ \text{is stable}\bigr).
\]

\subsection{Y.5.2 Quine-Stable Planners}
A planner~$\mathcal{P}$ is \emph{quine-stable} if each refinement step yields a
media quine.  
We prove that
\[
\Box_{\mathrm{Polyx}}\bigl(\mathcal{P}\ \text{is task-safe}\bigr).
\]

\subsection{Y.5.3 Verified Controllers}
A controller is \emph{verified} if it satisfies both modal predicates.  
We show that verified controllers converge in finite time to harmonic
normal forms in the collective manifold.


\section{Full Mathematical Formalisation of Semantic Robotics}
\label{sec:Y6}

We now summarise the algebraic and geometric structures that enable a full
formalisation of robotics within the RSVP--Polyxan framework.

\subsection{Y.6.1 Actions as Field Deformations}
Each control action~$a$ is interpreted as a compactly supported deformation:
\[
\mathcal{A}(a) : (\Phi,\mathbf{v},S) \to (\Phi',\mathbf{v}',S').
\]

\subsection{Y.6.2 Policies as Polyxan Rewrite Systems}
A policy~$\pi$ is encoded as a Polyxan rewrite system with curvature and
entropy tags.  
Policy refinement corresponds to EHR rewriting.

\subsection{Y.6.3 Verification via Synthetic Duality}
Using the duality functor~$\mathbb{D}$, the continuous safety properties of
actions correspond to discrete quine-stability conditions of plans.

\subsection{Y.6.4 The Robotics Collapse Theorem}
Finally, we establish the collapse analogue for robotics:

\begin{theorem}[Robotics Collapse Theorem]
Every verified robotic policy trained on semantic proxies collapses in finite
time to a unique harmonic-quine-stable controller.  
This controller is the curvature-minimising reconciliation of all observed
affordance models.
\end{theorem}

\begin{quote}
\textbf{Embodied agents, like semantic galaxies, converge.}  
When their models are curvature-matched and quine-stable,  
their behaviour becomes inevitable, predictable, and safe.
\end{quote}
